\documentclass[10pt,a4paper]{amsart}
\usepackage[margin=19mm]{geometry}
\usepackage{amsmath,amssymb,mathtools}
\usepackage[scr=rsfs,cal=euler]{mathalfa}
\usepackage{xfrac}
\usepackage[unicode,hidelinks,bookmarks=false]{hyperref}
\newcommand{\ie}{i.e.\ }
\newcommand{\cf}{cf.\ }
\newcommand{\resp}{respectively\ }
\newcommand{\Subsection}{\subsection}
\providecommand{\nobreakdash}{\nobreak\hbox{-}}
\setlength{\emergencystretch}{2em}
\multlinegap0pt
\usepackage{amssymb}
\usepackage{mathtools}
\newtheorem{lem}{Lemma}
\newcommand{\Ga}{\Gamma}
\DeclareMathOperator{\Prob}{Prob}
\newcommand{\acts}{\curvearrowright}
\newcommand{\cM}{\mathcal{M}}
\newcommand{\cR}{\mathcal{R}}
\newcommand{\ve}{\varepsilon}
\title{Uniform amenability at infinity}
\author{Narutaka Ozawa}
\date{Source: arXiv:2604.22412v4 (CC BY 4.0)}
\begin{document}
\maketitle

\noindent\emph{Source and licence.} Uniform amenability at infinity. Version arXiv:2604.22412v4 (CC BY 4.0), distributed by arXiv under the Creative Commons Attribution 4.0 International licence, http://creativecommons.org/licenses/by/4.0/, as recorded in the arXiv licence metadata for this exact version and shown on the abstract page https://arxiv.org/abs/2604.22412v4. Full licence notice: \texttt{licenses/arxiv-2604.22412-CC-BY-4.0.txt}.\par\smallskip

\noindent\textit{Editorial scope.} Counted unit: the article's lem lem:borelamencriterion (source0.tex lines 600--608). The article's complete original proof of it is delivered with it (source0.tex lines 609--659, one \texttt{proof} environment of 51 lines). No mathematics is written, repaired or re-derived: the statement and the proof are the article's own lines, in the article's own notation, and no retained line is substituted or re-flowed.  Retained uncounted as the source-local context the proof needs: lines 559--567.  Nothing was omitted from any retained span, so the package carries no omission record.  No retained line contains a citation, so the wrapper carries no bibliography and no bibliography entry.  No retained line contains a figure environment, an image-inclusion command or drawing code, so no image is delivered and the package declares no asset dependency.  Editorial formatting only, in addition: an AMS \texttt{amsart} preamble that restates the article's own declarations and macros used by the retained lines, namely the environments \texttt{lem} on separate counters, together with the standard packages those lines need.  The retained environments print in this document as Lemma 1 in order of appearance, whereas the article prints the same items with the numbering of its own full document.  The complete native source of the article as distributed in the arXiv e-print for this exact version is bundled unchanged as \texttt{sources/199090/source0.tex}.
\begin{lem}\label{lem:fd} 
Let $\Ga$ be a countable group and 
$Z$ be a Borel $\Ga$-space. 
Suppose $\Delta\le\Ga$ is a subgroup and $Y\subset Z$ is 
a $\Delta$-invariant Borel subset that satisfy 
$gY\cap Y=\emptyset$ for $g\in\Ga\setminus\Delta$ and  
$\Ga Y=Z$.  
If $Y$ is amenable as a Borel $\Delta$-space, then $Z$ is amenable. 
\end{lem}

\begin{lem}\label{lem:borelamencriterion}
Let $\Ga$ be a countable group and 
$Z$ be a Borel $\Ga$-space. 
Suppose that the countable Borel equivalence 
relation $\cR_{\Ga\acts Z}$ is $\cM$-amenable 
and that $\{\Ga^x : x\in Z\}$ consists of 
a countable family of amenable subgroups. 
Then the Borel $\Ga$-space $Z$ is amenable.
\end{lem}

\begin{proof}
We may assume by a countable decomposition 
that all $\Ga^x$ are conjugate to a single subgroup $\Delta\le\Ga$. 
The Borel subset $Y \coloneq \{ x\in Z : \Ga^x =\Delta \}$ 
is invariant under the normalizer subgroup $N_\Ga(\Delta)$ 
of $\Delta$ in $\Ga$.
Thanks to Lemma~\ref{lem:fd}, 
it suffices to show that $Y$ is amenable as 
a Borel $N_\Ga(\Delta)$-space.
Thus, we may assume that $Y=Z$ and $N_\Ga(\Delta)=\Ga$, 
i.e., $\Delta = \ker(\Ga\acts Z)$ 
and the action $\bar{\Ga} \coloneq \Ga/\Delta\acts Z$ is free. 
It follows that $Z$ is amenable as a Borel $\bar{\Ga}$-space. 

As in Proof of Lemma~\ref{lem:fd},  
fix a lift $\varsigma\colon \bar{\Ga}\to\Ga$ and the cocycle
$\alpha\colon \Ga \times \bar{\Ga} \to \Delta$.
Thus the bijection
$\theta\colon \bar{\Ga} \times \Delta \to\Ga$, given by 
 $(p,t) \mapsto \varsigma(p)t$, 
satisfies $s\theta(p,t) = \theta( sp,\alpha(s,p)t )$. 
For $\eta\colon Z\to \Prob \bar{\Ga}$ 
and $\xi\in\Prob \Delta$, we define 
$\zeta\colon Z\to\Prob \Ga$ by 
$\zeta^x = (\eta^x\otimes \xi)\circ\theta^{-1}$. 
Then 
\begin{align*}
\| \zeta^{sx}-s\zeta^x\|_1
 &\le \| \zeta^{sx}-(s\eta^x\otimes\xi)\circ\theta^{-1}\|_1
   + \| (s\eta^x\otimes\xi)\circ\theta^{-1}- s\zeta^x\|_1\\
 &\le \| \eta^{sx}-s\eta^x\|_1 
   + \sum_p\eta^x(s^{-1}p)\|\xi-\alpha(s,s^{-1}p)\xi\|_1.
\end{align*}
Let a finite subset $F\Subset\Ga$, $m\in\Prob Z$, and $\ve>0$ be given. 
There is $\eta$ that satisfies 
\[
\sum_{s\in F}\int_Z \| \eta^{sx}-s\eta^x\|_1\,dm(x) < \ve/2.
\]
Since $\sum_{s,p}\int_Z \eta^x(s^{-1}p)\,dm(x) = |F|$, 
there is $\xi$ that satisfies 
\[
\sum_{s\in F}\int_Z \sum_p\eta^x(s^{-1}p)\|\xi-\alpha(s,s^{-1}p)\xi\|_1\,dm(x) <\ve/2.
\]
It follows that for these $\eta$ and $\xi$, 
the Borel map $\zeta\colon Z\to\Prob \Ga$ satisfies\\ 
\begin{minipage}{.9\textwidth}
\[
\sum_{s\in F}\int_Z \| \zeta^{sx}-s\zeta^x\|_1\,dm(x) < \ve.
\]
\end{minipage} 
\end{proof}

\end{document}
